% ============================================================
%  sec-5-project-motivation.tex  —  1.5 Project Motivation
% ============================================================

\section{Project Motivation}
\label{sec:project-motivation}

The proliferation of generative AI and large-scale machine learning workloads
has fundamentally transformed modern computing infrastructure. Organisations
worldwide are investing heavily in AI-capable data centres, yet the prevailing
reliance on \acrshort{gpu}-dominated infrastructures has revealed critical
limitations in cost, energy efficiency, supply chain resilience, and strategic
autonomy. This project addresses the urgent industrial need for structured
methodologies to evaluate, guide, and optimise the migration from GPU-centric
architectures toward \acrshort{npu}-based alternatives.

% ── 1.5.1 ────────────────────────────────────────────────────
\subsection{Need for GPU\texorpdfstring{$\,\to\,$}{-to-}NPU Migration Guidance}
\label{subsec:migration-need}

Since the advent of deep learning, \acrfull{gpu}s have established themselves
as the de facto hardware foundation for AI workloads, with the market projected
to reach USD 228.04B by 2030.

\subsubsection{Structural Limitations Driving Migration}

Despite this dominance, the GPU-centric paradigm is increasingly
unsustainable. \tabref{gpu-limitations} summarises the five structural
limitations that motivate the exploration of NPU-based alternatives.

\begin{table}[htbp]
  \centering
  \caption{Structural limitations of GPU-centric AI infrastructure}
  \label{tab:gpu-limitations}
  \begin{tabularx}{\textwidth}{lX}
    \toprule
    \tableheader{Limitation & Key Evidence} \\
    \midrule
    Energy consumption
      & A single GB200 NVL72 rack draws 120–132 kW vs.\ 10–15 kW for a typical
        enterprise rack; AI's energy needs projected at 21\% of global
        electricity by 2030. \\
    \rowcolor{TableRowOdd}
    Infrastructure cost
      & Data centre construction costs rose from \$7.7M to \$10.7M per MW
        between 2020 and 2025 (7\% CAGR). \\
    Vendor lock-in
      & CUDA's dominance creates prohibitive migration barriers; multi-year supply
        agreements between NVIDIA and hyperscalers crowd out enterprise buyers. \\
    \rowcolor{TableRowOdd}
    Scalability for mid-size enterprises
      & Four GB200 racks require $\approx$0.5 MW of dedicated power plus
        liquid cooling infrastructure unavailable in most enterprise data centres. \\
    Supply chain constraints
      & GPU delivery windows stretch 36–52 weeks; Blackwell-class hardware
        delayed to Q1 2027; NVIDIA consumes $\approx$60\% of CoWoS packaging
        capacity. \\
    \bottomrule
  \end{tabularx}
\end{table}

NPUs address these limitations through domain-specific architectures.
Recent benchmarks show NPUs delivering 30–63\% higher energy efficiency than
NVIDIA GPUs in quantised \acrshort{llm} inference, while operating at
10–30W per chip compared to 300–700W for data centre GPUs.

\subsubsection{Enterprise Migration Challenges}

Despite the promise of NPUs, enterprises face significant migration barriers:
GPU-oriented ecosystems, unclear migration paths, limited NPU expertise, and
workload heterogeneity (training vs.\ inference, latency vs.\ throughput, cloud
vs.\ edge). The absence of structured guidance creates a systematic gap between
the theoretical advantages of NPU adoption and practical deployment capability.
This project directly addresses this gap.

% ── 1.5.2 ────────────────────────────────────────────────────
\subsection{Need for Compatibility Analysis}
\label{subsec:compatibility-need}

Compatibility analysis is a critical technical prerequisite during AI hardware
transitions. The shift from GPU to NPU is not merely a hardware substitution
but a fundamental reconfiguration of the entire software-hardware stack. A 2025
industry analysis found that 68\% of production AI workloads contain at least
one CUDA-dependent component requiring manual reimplementation for NPU
deployment, rising to 85\% for workloads using custom architectures.

\subsubsection{Software Compatibility Dimensions}

\begin{table}[htbp]
  \centering
  \caption{AI framework compatibility status across hardware backends}
  \label{tab:framework-compat}
  \begin{tabularx}{\textwidth}{llXl}
    \toprule
    \tableheader{Framework & Primary Backend & NPU Support Status & Migration Complexity} \\
    \midrule
    PyTorch      & CUDA (NVIDIA)    & Limited (requires CANN/MindSpore adaptation) & High \\
    \rowcolor{TableRowOdd}
    TensorFlow   & CUDA (NVIDIA)    & Partial (via Ascend TF plugin)               & Medium-High \\
    MindSpore    & Ascend NPU       & Native                                       & Low \\
    \rowcolor{TableRowOdd}
    ONNX Runtime & Multi-arch       & Moderate (provider-dependent)                & Medium \\
    \bottomrule
  \end{tabularx}
\end{table}

CUDA creates dependency at multiple abstraction levels simultaneously:
custom CUDA kernels, cuDNN/cuBLAS libraries, profiling tools (Nsight), and
model checkpoint formats with CUDA-specific metadata. Each layer must be
validated independently during migration.

\subsubsection{Impact of Compatibility Problems}

Compatibility failures affect four dimensions of deployment success.
\textbf{Feasibility}: severe incompatibilities can render deployment impossible
without prohibitive re-engineering — cases include models with proprietary
CUDA kernels and framework versions incompatible with NPU drivers.
\textbf{Performance}: suboptimal operator implementations may reduce throughput
by 30–50\% and increase latency through memory layout mismatches.
\textbf{Maintenance complexity}: organisations maintaining both GPU and NPU
infrastructure must support two complete, evolving software stacks.
\textbf{Cost}: industry estimates suggest 40–60\% of total migration cost
arises from software compatibility resolution, not hardware procurement.

Such migrations require not only code changes but framework-level retraining,
model checkpoint conversion, and accuracy validation — illustrating the need
for systematic compatibility analysis methodologies.

% ── 1.5.3 ────────────────────────────────────────────────────
\subsection{Need for Hardware Transition Evaluation}
\label{subsec:evaluation-need}

Hardware transition decisions represent multi-year infrastructure commitments
that shape organisational AI capabilities. Systematic evaluation must span
five complementary dimensions.

\subsubsection{Performance Benchmarking}

GPU and NPU performance profiles differ fundamentally across workload types.
For inference at batch size 1, NPUs can achieve 3–10$\times$ lower latency
than GPUs. For large-scale training, GPUs remain significantly superior.
Benchmarks must be conducted under realistic production conditions — including
data pipeline overhead, network latency, and memory contention — not idealised
laboratory settings.

This project addresses the need for systematic methodologies, enabling
organisations to make data-driven infrastructure decisions based on
comprehensive technical and economic assessment rather than incomplete
information or vendor marketing claims.
